% ============================================================================
\section{Scaling Bitwise Computation to Large Language Models}
\label{sec:scaling-llms}
% ============================================================================

Armed with an understanding of the transformer architecture, we can
now revisit the question from \cref{sec:hybrid}: can the four bitwise
paradigms be combined in a way that scales to LLM-sized models?

The previous section showed that only two components---the attention
softmax and layer normalisation---fundamentally require floating-point
arithmetic. This section proposes a combination strategy that assigns
each paradigm to the transformer component it is best suited for.

% ----------------------------------------------------------------------------
\subsection{Why No Single Paradigm Scales}
\label{sec:why-no-single}
% ----------------------------------------------------------------------------

Let us first be clear about why each paradigm \emph{alone} fails at
LLM scale:

\begin{description}[leftmargin=1em]
  \item[Binary Neural Networks] can binarise weights and activations,
    but the BNN approach assumes a fixed network topology. It does not
    address the attention mechanism, which is the defining innovation
    of transformers. BitNet b1.58~\cite{ma2024bitnet} succeeds
    precisely because it only binarises the \emph{weights}---not the
    architecture.

  \item[Tsetlin Machines] learn propositional clauses over Boolean
    features. Language tokens are categorical with distributional
    semantics---booleanising 4{,}096-dimensional embeddings would
    destroy the fine-grained similarity structure that makes language
    modelling work. Additionally, Tsetlin Machines have no native
    mechanism for the pairwise token interactions that attention
    provides.

  \item[Logic Gate Networks] with 2-input gates would require
    astronomical circuit depth to approximate the $n \times n$
    attention computation over thousands of tokens. The softmax
    operation is inherently continuous. However, LGNs could work at
    the \emph{component level}---replacing individual MLP blocks
    rather than the entire transformer.

  \item[KAN / LUT compilation] faces an exponential scaling wall:
    lookup tables grow exponentially with input dimensionality.
    A function of 4{,}096-dimensional embeddings cannot be stored
    as a table. But a \emph{univariate} activation function (1D input)
    can easily be stored as a 256-entry table, which is exactly where
    LUT compilation shines.
\end{description}

\begin{keyinsight}
The insight is not ``which paradigm replaces the transformer?'' but
rather ``which paradigm can serve as a component \emph{inside} a
transformer?'' The transformer architecture is non-negotiable for
language. The four paradigms become optimisation techniques within
that structure, not replacements for it.
\end{keyinsight}

% ----------------------------------------------------------------------------
\subsection{A Combined Architecture: Bitwise Components in a Transformer Shell}
\label{sec:combined-architecture}
% ----------------------------------------------------------------------------

We propose mapping each paradigm to the transformer component where
its strengths are most relevant:

\begin{figure}[htbp]
\centering
\begin{tikzpicture}[
  comp/.style={draw, rounded corners=4pt, minimum width=5cm,
    minimum height=1.0cm, font=\small\sffamily, align=center,
    line width=0.8pt},
  paradigm/.style={font=\scriptsize\sffamily\bfseries, rounded corners=2pt,
    draw, inner sep=3pt, fill=#1!15, draw=#1!60},
  arr/.style={-{Stealth[length=5pt]}, thick},
  note/.style={font=\scriptsize\sffamily, text width=5cm}
]
  % Components stack
  \node[comp, fill=bitgray!5] (emb) at (0, 0)
    {Embedding lookup};
  \node[comp, fill=bitblue!8] (qkv) at (0, 1.5)
    {QKV projection};
  \node[comp, fill=red!8, draw=red!40] (attn) at (0, 3)
    {Attention softmax};
  \node[comp, fill=bitblue!8] (oproj) at (0, 4.5)
    {Output projection};
  \node[comp, fill=red!8, draw=red!40] (ln) at (0, 6)
    {Layer normalisation};
  \node[comp, fill=bitorange!8] (mlp) at (0, 7.5)
    {MLP / feed-forward};
  \node[comp, fill=bitgreen!8] (gate) at (0, 9)
    {MoE router (if used)};

  % Paradigm labels
  \node[paradigm=bitgray, right] at (3.5, 0)
    {Integer table lookup (no paradigm needed)};
  \node[paradigm=bitblue, right] at (3.5, 1.5)
    {BitNet: ternary $\{-1, 0, +1\}$ weights};
  \node[paradigm=red, right] at (3.5, 3)
    {Open problem (see \cref{sec:softmax-problem})};
  \node[paradigm=bitblue, right] at (3.5, 4.5)
    {BitNet: ternary weights};
  \node[paradigm=red, right] at (3.5, 6)
    {Open problem (see \cref{sec:layernorm-problem})};
  \node[paradigm=bitorange, right] at (3.5, 7.5)
    {LGN inference + LUT activations};
  \node[paradigm=bitgreen, right] at (3.5, 9)
    {Tsetlin Machine gating};

  % Arrows
  \foreach \a/\b in {emb/qkv, qkv/attn, attn/oproj, oproj/ln, ln/mlp, mlp/gate} {
    \draw[arr] (\a) -- (\b);
  }

  % Legend
  \draw[thick, red!40, fill=red!8, rounded corners=2pt]
    (-3.5, -1.5) rectangle (-1.5, -1.0);
  \node[font=\scriptsize\sffamily, right] at (-1.3, -1.25)
    {= remaining floating-point};
  \draw[thick, bitblue!40, fill=bitblue!8, rounded corners=2pt]
    (-3.5, -2.2) rectangle (-1.5, -1.7);
  \node[font=\scriptsize\sffamily, right] at (-1.3, -1.95)
    {= integer/ternary arithmetic};
\end{tikzpicture}
\caption{Proposed mapping of bitwise paradigms to transformer
  components. Green and blue components can be made (near-)bitwise
  today. Red components remain open problems requiring floating-point
  arithmetic or novel replacements.}
\label{fig:combined-architecture}
\end{figure}

\subsubsection{Component 1: BitNet Ternary Weights (BNN Paradigm)}

The weight matrices $\bm{W}_Q$, $\bm{W}_K$, $\bm{W}_V$,
$\bm{W}_O$, $\bm{W}_1$, $\bm{W}_2$ are constrained to ternary
values $\{-1, 0, +1\}$, following the BitNet b1.58
approach~\cite{ma2024bitnet}.

\textbf{What this buys:}
\begin{itemize}[leftmargin=2em]
  \item Weight storage drops from 16~bits to 1.58~bits per parameter
    (a 10$\times$ reduction).
  \item Matrix multiplication becomes addition and subtraction (no
    multiplier circuits needed).
  \item Proven at 3B+ parameter scale with competitive accuracy.
\end{itemize}

\textbf{What remains floating-point:} The activations (intermediate
values flowing between layers) remain in higher precision (typically
8-bit integers in BitNet). The ternary constraint applies only to the
stored weights.

\subsubsection{Component 2: LUT-Compiled Activations (KAN Paradigm)}

The activation functions in the MLP blocks (GELU, SiLU, etc.)
are fixed mathematical functions applied element-wise. Instead of
computing them with floating-point arithmetic, we can:

\begin{enumerate}[leftmargin=2em]
  \item \textbf{Train} with learnable activation functions
    (KAN-style B-splines or learnable piecewise-linear functions).
  \item \textbf{Compile} the learned functions into lookup tables:
    a 256-entry table maps each possible 8-bit input to a
    precomputed output.
\end{enumerate}

\textbf{What this buys:}
\begin{itemize}[leftmargin=2em]
  \item Activation computation becomes a single table lookup---pure
    memory access, no arithmetic.
  \item Each layer can learn its own optimal activation shape, rather
    than using one fixed function everywhere.
  \item The LUT fits in L1 cache (256 entries $\times$ 1~byte = 256~bytes).
\end{itemize}

\textbf{Why KANs do not replace the full MLP:} As discussed in
\cref{sec:ffn}, the MLP block has $d \times d_\text{ff}$ connections.
Replacing each with a full B-spline would create an explosion in
parameters and training cost. But replacing only the
\emph{per-neuron activation function} with a learned spline and then
compiling to a LUT is lightweight and practical.

\subsubsection{Component 3: Logic Gate Network MLP (LGN Paradigm)}

The MLP blocks account for roughly two-thirds of the parameters and
computation in a transformer. After training with ternary weights and
LUT activations, the MLP can potentially be \emph{further compiled}
into a logic gate network for inference:

\begin{enumerate}[leftmargin=2em]
  \item Train the MLP with standard methods (ternary weights + LUT
    activations).
  \item After training, \emph{compile} the trained MLP into an
    equivalent Boolean circuit, where the quantised weights and
    activations are mapped to logic gate configurations.
  \item The inference-time MLP becomes a fixed logic circuit with no
    arithmetic operations.
\end{enumerate}

\textbf{What this buys:}
\begin{itemize}[leftmargin=2em]
  \item Inference-time MLP is pure combinational logic (AND/OR/XOR
    gates)---maximum hardware efficiency.
  \item Maps directly to FPGA or ASIC implementations.
  \item No weight memory needed at inference: the function is
    ``burned into'' the circuit topology.
\end{itemize}

\textbf{What is unproven:} LGNs have been demonstrated at MNIST scale
(thousands of gates). A transformer MLP block at 7B scale has 67M
parameters per layer. Whether the LGN compilation approach can scale
to this size without prohibitive circuit depth or accuracy loss is an
open research question.

\subsubsection{Component 4: Tsetlin Machine Router (TM Paradigm)}

In a Mixture-of-Experts transformer, the gating network decides which
expert processes each token. This is a small classification problem:
given a token's hidden state, select the top-$k$ experts from a pool
of $E$ options.

Tsetlin Machines are a natural fit here because:
\begin{itemize}[leftmargin=2em]
  \item The routing decision is a low-dimensional classification
    (typically $E = 8$ to 64 experts).
  \item The input can be booleanised (e.g., by thresholding the
    hidden state dimensions) without catastrophic information loss,
    because the routing decision does not need the same fine-grained
    precision as the main computation.
  \item The learned clauses would be \emph{interpretable}: we could
    inspect \emph{why} a particular token was routed to a particular
    expert, expressed as human-readable Boolean rules.
\end{itemize}

\textbf{What this buys:}
\begin{itemize}[leftmargin=2em]
  \item Interpretable routing decisions (no black-box gating).
  \item Integer-only routing computation (no floating-point softmax
    in the gating network).
  \item Potential for faster routing than a learned linear layer.
\end{itemize}

\textbf{What is unproven:} Whether Tsetlin Machine routing can match
the load-balancing and accuracy of continuous gating functions at LLM
scale.

% ----------------------------------------------------------------------------
\subsection{The Softmax Problem}
\label{sec:softmax-problem}
% ----------------------------------------------------------------------------

The attention softmax (\cref{eq:softmax}) is the single most
important obstacle to a fully bitwise transformer. Let us examine why
it is hard and what alternatives exist.

\subsubsection{Why Softmax Resists Binarisation}

Softmax converts a vector of raw scores $\bm{z} = [z_1, \ldots, z_n]$
into a probability distribution:

\[
  \text{softmax}(z_i) = \frac{e^{z_i}}{\sum_{j=1}^{n} e^{z_j}}
\]

This operation has three properties that are fundamentally at odds
with bitwise computation:

\begin{enumerate}[leftmargin=2em]
  \item \textbf{Exponentiation.} The function $e^x$ maps integers to
    irrational numbers. There is no bitwise approximation of $e^x$
    that preserves its essential property: amplifying differences
    between large and small values.

  \item \textbf{Normalisation.} Dividing by the sum ensures outputs
    sum to~1. Division is the most expensive basic arithmetic
    operation and does not have a bitwise equivalent.

  \item \textbf{Precision sensitivity.} The attention mechanism
    relies on distinguishing between scores like 0.15 and 0.17. In
    a 4{,}096-token sequence, the softmax distributes attention mass
    across thousands of tokens, and small differences in the
    distribution affect which information propagates forward. Binary
    attention (attend or not) would be too coarse.
\end{enumerate}

\subsubsection{Potential Replacements}

Several research directions aim to reduce or eliminate the softmax:

\begin{description}[leftmargin=1em]
  \item[Linear attention.] Replace $\text{softmax}(\bm{Q}\bm{K}^\top)\bm{V}$
    with $\phi(\bm{Q})\phi(\bm{K})^\top\bm{V}$, where $\phi$ is a
    feature map (e.g., $\text{elu}(x) + 1$ or random
    features). This avoids the $n \times n$ attention matrix entirely.
    Linear attention has $O(n)$ complexity instead of $O(n^2)$, but
    current models using it (e.g., RWKV, Mamba) trade some quality
    for this efficiency.

  \item[ReLU or squared attention.] Replace softmax with
    $\text{ReLU}(\bm{Q}\bm{K}^\top)$ or
    $(\bm{Q}\bm{K}^\top)^2$. These are non-negative (like softmax)
    but avoid exponentiation. Squaring is just multiplication;
    ReLU is a simple threshold (which \emph{is} a bitwise-friendly
    operation). Recent work shows these can achieve competitive
    quality on some tasks, but they lack softmax's normalisation,
    which can cause training instability.

  \item[Integer softmax approximation.] Approximate $e^x$ using
    piecewise linear functions or lookup tables, and approximate
    division using bit shifts. This has been demonstrated in
    inference-only quantisation (e.g., I-BERT), but typically with
    some accuracy degradation.

  \item[Gated attention / state-space models.] Architectures like
    Mamba~\cite{gu2024mamba} and RWKV replace attention entirely
    with recurrent state-space mechanisms that use element-wise
    gating instead of softmax. These are more amenable to
    integer/bitwise implementation, and have shown competitive
    results on language modelling benchmarks at multi-billion
    parameter scale.
\end{description}

\begin{keyinsight}
The softmax problem may not need to be ``solved'' within the
softmax framework. If state-space models (like Mamba) or linear
attention variants prove competitive with transformers at scale,
they offer a path to fully integer attention. The question shifts
from ``how do we binarise softmax?'' to ``can we build LLMs without
softmax at all?'' Early results are promising but not yet
conclusive at the largest scales.
\end{keyinsight}

% ----------------------------------------------------------------------------
\subsection{The Layer Normalisation Problem}
\label{sec:layernorm-problem}
% ----------------------------------------------------------------------------

Layer normalisation (\cref{eq:layernorm}) is the second
floating-point holdout. It requires computing the mean and variance
of a vector, then dividing each element by the standard deviation.

The situation here is more tractable than softmax:

\begin{description}[leftmargin=1em]
  \item[RMSNorm] (used in Llama and most modern LLMs) drops the
    mean computation entirely:
    $\text{RMSNorm}(\bm{x}) = \bm{x} / \text{RMS}(\bm{x}) \cdot \bm{\gamma}$,
    where $\text{RMS}(\bm{x}) = \sqrt{\frac{1}{d}\sum_i x_i^2}$.
    This still requires a square root and division, but the
    computation is simpler.

  \item[Integer RMSNorm.] The sum of squares $\sum_i x_i^2$ is an
    integer operation if $\bm{x}$ is integer. The square root and
    division can be approximated with bit-shift techniques or small
    lookup tables. Practical integer-only implementations have been
    demonstrated in quantised inference.

  \item[BitNet's approach.] BitNet b1.58 retains RMSNorm in
    higher precision, arguing that the normalisation computation is
    a tiny fraction of total operations (it processes $d$ values
    per token, versus $d \times d_\text{ff}$ in the MLP). The
    pragmatic argument: even if normalisation stays in floating
    point, it accounts for less than 1\% of total operations.
\end{description}

\begin{engineernote}
Layer normalisation may be the wrong hill to die on for full
binarisation. Even in a hypothetical ``99\% bitwise'' transformer,
the normalisation steps contribute a negligible fraction of total
energy and latency. The engineering priority should be the
attention mechanism and MLP blocks, which dominate both computation
and memory bandwidth.
\end{engineernote}

% ----------------------------------------------------------------------------
\subsection{The Realistic Picture}
\label{sec:realistic-picture}
% ----------------------------------------------------------------------------

\Cref{tab:bitwise-llm-summary} summarises the current state of each
component.

\begin{table}[htbp]
\centering
\caption{Feasibility assessment for bitwise transformer components.
  ``Proven'' means demonstrated at LLM scale (billions of parameters);
  ``plausible'' means demonstrated at smaller scale or with minor
  accuracy trade-offs; ``open'' means no convincing demonstration
  exists.}
\label{tab:bitwise-llm-summary}
\begin{tabular}{@{}p{3cm}lp{3cm}p{4.5cm}@{}}
  \toprule
  \textbf{Component} & \textbf{Paradigm} & \textbf{Status}
    & \textbf{Key challenge} \\
  \midrule
  Weight storage & BitNet (BNN) & Proven (3B+) & None---ternary weights work. \\
  Matrix multiply & BitNet (BNN) & Proven & Addition/subtraction only. \\
  Activation fn. & KAN $\to$ LUT & Plausible & Per-layer learned activations need validation at scale. \\
  MLP compilation & LGN & Open & Circuit depth at transformer MLP scale. \\
  MoE routing & Tsetlin Machine & Open & Load balancing, accuracy matching. \\
  Attention & (see below) & Open & Softmax replacement. \\
  Layer norm & Integer approx. & Plausible & Small accuracy trade-off acceptable. \\
  \bottomrule
\end{tabular}
\end{table}

\subsubsection{What Could Work Today}

A ``mostly bitwise'' LLM using:
\begin{itemize}[leftmargin=2em]
  \item Ternary weights throughout (BitNet b1.58),
  \item LUT-compiled activation functions,
  \item Integer-approximated RMSNorm,
  \item Standard softmax attention (the remaining floating-point
    component).
\end{itemize}

This configuration would make approximately 95\% of operations
integer or bitwise, with floating point confined to the attention
softmax and a thin normalisation layer. It is achievable with
today's research.

\subsubsection{What Could Work Soon}

Replace softmax attention with a softmax-free alternative:
\begin{itemize}[leftmargin=2em]
  \item State-space models (Mamba) or linear attention for the
    sequence mixing component.
  \item ReLU attention or gated attention as drop-in replacements
    within the transformer structure.
  \item Integer-approximated softmax via LUT-based $e^x$ and
    bit-shift division.
\end{itemize}

This would push the architecture to 99\%+ bitwise/integer operations.
The research exists in pieces; the integration and validation at
scale is the remaining work.

\subsubsection{What Remains Speculative}

A \emph{fully} bitwise LLM with:
\begin{itemize}[leftmargin=2em]
  \item Logic Gate Network MLP blocks (compiled from trained ternary
    MLPs),
  \item Tsetlin Machine MoE routing,
  \item All-integer attention and normalisation.
\end{itemize}

This is the ambitious end-state. Each component exists in isolation
at small scale, but the integration, the compounding of
approximation errors across 80+ layers, and the training dynamics
at billions of parameters are all uncharted territory.

% ----------------------------------------------------------------------------
\subsection{The Path Forward}
\label{sec:path-forward-llms}
% ----------------------------------------------------------------------------

Rather than attempting a fully bitwise LLM immediately, a staged
research programme would build confidence incrementally:

\begin{enumerate}[leftmargin=2em]
  \item \textbf{MNIST hybrid} (current project): Validate the
    four-paradigm combination on a tractable problem.

  \item \textbf{BitNet + LUT activations}: Take an existing BitNet
    model and replace fixed activations with LUT-compiled learned
    activations. Measure accuracy impact. This is a minimal
    intervention with clear evaluation.

  \item \textbf{Softmax-free BitNet}: Combine ternary weights with
    a softmax-free attention mechanism (linear attention, ReLU
    attention, or Mamba-style gating). Evaluate on standard language
    benchmarks.

  \item \textbf{LGN MLP compilation}: Train a small transformer
    (125M parameters) with ternary weights, then attempt to compile
    the MLP blocks into logic gate networks. Measure inference
    speed and accuracy preservation.

  \item \textbf{TM routing for MoE}: Implement Tsetlin Machine
    routing in a MoE transformer and compare with standard softmax
    gating on load balance and accuracy.

  \item \textbf{Integration}: Combine the successful components
    into a single architecture and evaluate at increasing scale.
\end{enumerate}

Each step is independently publishable and valuable regardless of
whether the full vision materialises. The MNIST project we develop
in this article is step~1.
